\section*{References}\label{references}

\begin{enumerate}
\def\labelenumi{\arabic{enumi}.}
\item
  Kamzolas, I. et al.~A data-driven framework reconstructs the molecular continuum of human MASLD progression. \emph{Nat Metab} \textbf{8}, 1545--1562 (2026). \url{https://doi.org/10.1038/s42255-026-01543-7}
\item
  Rivas, M. A. et al.~Effect of predicted protein-truncating genetic variants on the human transcriptome. \emph{Science} \textbf{348}, 666--669 (2015). \url{https://doi.org/10.1126/science.1261877}
\item
  Kendall, T. J. et al.~An integrated gene-to-outcome multimodal database for metabolic dysfunction-associated steatotic liver disease. \emph{Nat Med} \textbf{29}, 2939--2953 (2023). \url{https://doi.org/10.1038/s41591-023-02602-2}
\item
  Alvarez-Sola, G. et al.~LiRNA: An interactive atlas of human liver RNAseq databases. \emph{JHEP Rep} \textbf{8}, 101914 (2026). \url{https://doi.org/10.1016/j.jhepr.2026.101914}
\item
  Piras, I. S. \& DiStefano, J. K. Comprehensive meta-analysis reveals distinct gene expression signatures of MASLD progression. \emph{Life Sci Alliance} \textbf{7}, e202302517 (2024). \url{https://doi.org/10.26508/lsa.202302517}
\item
  Gobeil, É. et al.~Integrative genetic and liver transcriptomic analyses identify TRIB1AL as a target for steatotic liver disease. \emph{J Clin Endocrinol Metab} \textbf{111}, 1980--1993 (2026a). \url{https://doi.org/10.1210/clinem/dgag043}
\item
  Saliba-Gustafsson, P. et al.~A functional genomic framework to elucidate novel causal metabolic dysfunction-associated fatty liver disease genes. \emph{Hepatology} \textbf{82}, 165--183 (2025). \url{https://doi.org/10.1097/HEP.0000000000001066}
\item
  Jamialahmadi, O. et al.~Partitioned polygenic risk scores identify distinct types of metabolic dysfunction-associated steatotic liver disease. \emph{Nat Med} \textbf{30}, 3614--3623 (2024). \url{https://doi.org/10.1038/s41591-024-03284-0}
\item
  Currin, K. W. et al.~Genetic effects on chromatin accessibility uncover mechanisms of liver gene regulation and quantitative traits. \emph{Genome Res} \textbf{35}, 1485--1502 (2025). \url{https://doi.org/10.1101/gr.279741.124}
\item
  Broadaway, K. A. et al.~Liver eQTL meta-analysis illuminates potential molecular mechanisms of cardiometabolic traits. \emph{Am J Hum Genet} \textbf{111}, 1899--1913 (2024). \url{https://doi.org/10.1016/j.ajhg.2024.07.017}
\item
  Ghodsian, N. et al.~Electronic health record-based genome-wide meta-analysis provides insights on the genetic architecture of non-alcoholic fatty liver disease. \emph{Cell Rep Med} \textbf{2}, 100437 (2021). \url{https://doi.org/10.1016/j.xcrm.2021.100437}
\item
  Ji, Y. et al.~Human liver methionine cycle: MAT1A and GNMT gene resequencing, functional genomics, and hepatic genotype-phenotype correlation. \emph{Drug Metab Dispos} \textbf{40}, 1984--1992 (2012). \url{https://doi.org/10.1124/dmd.112.046953}
\item
  Gobeil, É. et al.~An integrated proteogenomic investigation of the human liver uncovers molecular drivers of steatotic liver disease. Preprint at medRxiv \url{https://doi.org/10.64898/2026.06.04.26354903} (2026b).
\item
  Kumak, E., Darde, T. \& Konu, O. Network-based meta-analysis maps stage-dependent molecular programs in MASLD through MASLD-META NETWORK application. Preprint at bioRxiv \url{https://doi.org/10.64898/2026.08.26.747338} (2026).
\item
  Filliol, A. et al.~Opposing roles of hepatic stellate cell subpopulations in hepatocarcinogenesis. \emph{Nature} \textbf{610}, 356--365 (2022). \url{https://doi.org/10.1038/s41586-022-05289-6}
\item
  Guilliams, M. et al.~Spatial proteogenomics reveals distinct and evolutionarily conserved hepatic macrophage niches. \emph{Cell} \textbf{185}, 379--396.e38 (2022). \url{https://doi.org/10.1016/j.cell.2021.12.018}
\item
  Vu, H. et al.~Progressive fibrosis in human MASLD is associated with spatially linked transcriptomic signatures of metabolic reprogramming and senescence. \emph{JHEP Rep} \textbf{8}, 101657 (2026). \url{https://doi.org/10.1016/j.jhepr.2025.101657}
\item
  Boesch, M. et al.~Integrated multi-omics identifies distinct macrophage alterations during progression of metabolic dysfunction-associated steatohepatitis. \emph{Nat Genet} \textbf{58}, 1295--1307 (2026). \url{https://doi.org/10.1038/s41588-026-02600-3}
\item
  Yakubovsky, O. et al.~A spatial atlas of the healthy human liver from live donors. \emph{Nature} \textbf{653}, 1148--1157 (2026). \url{https://doi.org/10.1038/s41586-026-10377-y}
\item
  Weiss, C. A. M. et al.~Single-cell spatial proteomics maps human liver zonation patterns and their vulnerability to disruption in tissue architecture. \emph{Nat Metab} \textbf{8}, 741--756 (2026). \url{https://doi.org/10.1038/s42255-026-01459-2}
\item
  Hoang, S. A. et al.~Gene expression predicts histological severity and reveals distinct molecular profiles of nonalcoholic fatty liver disease. \emph{Sci Rep} \textbf{9}, 12541 (2019). \url{https://doi.org/10.1038/s41598-019-48746-5}
\item
  Govaere, O. et al.~Transcriptomic profiling across the nonalcoholic fatty liver disease spectrum reveals gene signatures for steatohepatitis and fibrosis. \emph{Sci Transl Med} \textbf{12}, eaba4448 (2020). \url{https://doi.org/10.1126/scitranslmed.aba4448}
\item
  Pantano, L. et al.~Molecular characterization and cell type composition deconvolution of fibrosis in NAFLD. \emph{Sci Rep} \textbf{11}, 18045 (2021). \url{https://doi.org/10.1038/s41598-021-96966-5}
\item
  Ramachandran, P. et al.~Resolving the fibrotic niche of human liver cirrhosis at single-cell level. \emph{Nature} \textbf{575}, 512--518 (2019). \url{https://doi.org/10.1038/s41586-019-1631-3}
\item
  Elison, W. et al.~Single cell multiomics reveals drivers of metabolic dysfunction-associated steatohepatitis. Preprint at medRxiv \url{https://doi.org/10.1101/2025.05.09.25327043} (2025).
\item
  Sveinbjornsson, G. et al.~Multiomics study of nonalcoholic fatty liver disease. \emph{Nat Genet} \textbf{54}, 1652--1663 (2022). \url{https://doi.org/10.1038/s41588-022-01199-5}
\item
  Oakes, C. \& Pachter, L. Transcriptomic responses to endurance exercise training in rats. \emph{BMC Genom Data} \textbf{27}, 74 (2026). \url{https://doi.org/10.1186/s12863-026-01440-8}
\item
  Gribben, C. et al.~Acquisition of epithelial plasticity in human chronic liver disease. \emph{Nature} \textbf{630}, 166--173 (2024). \url{https://doi.org/10.1038/s41586-024-07465-2}
\item
  Andrews, T. S. et al.~Single-cell, single-nucleus, and spatial RNA sequencing of the human liver identifies cholangiocyte and mesenchymal heterogeneity. \emph{Hepatol Commun} \textbf{6}, 821--840 (2022). \url{https://doi.org/10.1002/hep4.1854}
\item
  Namjou, B. et al.~GWAS and enrichment analyses of non-alcoholic fatty liver disease identify new trait-associated genes and pathways across eMERGE Network. \emph{BMC Med} \textbf{17}, 135 (2019). \url{https://doi.org/10.1186/s12916-019-1364-z}
\item
  Anstee, Q. M. et al.~Genome-wide association study of non-alcoholic fatty liver and steatohepatitis in a histologically characterised cohort. \emph{J Hepatol} \textbf{73}, 505--515 (2020). \url{https://doi.org/10.1016/j.jhep.2020.04.003}
\item
  Liu, Y. et al.~Genetic architecture of 11 organ traits derived from abdominal MRI using deep learning. \emph{eLife} \textbf{10}, e65554 (2021). \url{https://doi.org/10.7554/eLife.65554}
\item
  Haas, M. E. et al.~Machine learning enables new insights into genetic contributions to liver fat accumulation. \emph{Cell Genom} \textbf{1}, 100066 (2021). \url{https://doi.org/10.1016/j.xgen.2021.100066}
\item
  van der Meer, D. et al.~The link between liver fat and cardiometabolic diseases is highlighted by genome-wide association study of MRI-derived measures of body composition. \emph{Commun Biol} \textbf{5}, 1271 (2022). \url{https://doi.org/10.1038/s42003-022-04237-4}
\item
  Sinnott-Armstrong, N. et al.~Genetics of 35 blood and urine biomarkers in the UK Biobank. \emph{Nat Genet} \textbf{53}, 185--194 (2021). \url{https://doi.org/10.1038/s41588-020-00757-z}
\item
  Kanai, M. et al.~Genetic analysis of quantitative traits in the Japanese population links cell types to complex human diseases. \emph{Nat Genet} \textbf{50}, 390--400 (2018). \url{https://doi.org/10.1038/s41588-018-0047-6}
\item
  McCarthy, D. J. \& Smyth, G. K. Testing significance relative to a fold-change threshold is a TREAT. \emph{Bioinformatics} \textbf{25}, 765--771 (2009). \url{https://doi.org/10.1093/bioinformatics/btp053}
\item
  Pinheiro, J., Bates, D. \& R Core Team. nlme: Linear and Nonlinear Mixed Effects Models. R package version 3.1-168. \url{https://CRAN.R-project.org/package=nlme} (\url{doi:10.32614/CRAN.package.nlme}).
\item
  Gerard, D. Data-based RNA-seq simulations by binomial thinning. \emph{BMC Bioinformatics} \textbf{21}, 206 (2020). \url{https://doi.org/10.1186/s12859-020-3450-9}
\item
  Zhu, B. et al.~Integrative analyses elucidate transcriptional regulatory functions of risk alleles for metabolic liver disease. \emph{Nat Genet} \textbf{58}, 1280--1294 (2026). \url{https://doi.org/10.1038/s41588-026-02617-8}
\item
  Li, Z. et al.~Spatially resolved multi-omics of human metabolic dysfunction-associated steatotic liver disease. \emph{Nat Genet} \textbf{57}, 3112--3125 (2025). \url{https://doi.org/10.1038/s41588-025-02407-8}
\item
  Hong, S. E. et al.~Single-cell eQTL analysis identifies genetic variation underlying metabolic dysfunction-associated steatohepatitis. \emph{Nat Genet} \textbf{57}, 1638--1648 (2025). \url{https://doi.org/10.1038/s41588-025-02237-8}
\item
  Kim, H. Y. et al.~Multi-modal analysis of human hepatic stellate cells identifies novel therapeutic targets for metabolic dysfunction-associated steatotic liver disease. \emph{J Hepatol} \textbf{82}, 882--897 (2025). \url{https://doi.org/10.1016/j.jhep.2024.10.044}
\item
  Nasser, J. et al.~Genome-wide enhancer maps link risk variants to disease genes. \emph{Nature} \textbf{593}, 238--243 (2021). \url{https://doi.org/10.1038/s41586-021-03446-x}
\item
  Seagle, H. M. et al.~Genomics-informed drug-repurposing strategy identifies two therapeutic targets for preventing liver disease associated with metabolic dysfunction. \emph{Am J Hum Genet} \textbf{112}, 1778--1791 (2025). \url{https://doi.org/10.1016/j.ajhg.2025.06.014}
\item
  Kozumi, K. et al.~Transcriptomics identify thrombospondin-2 as a biomarker for NASH and advanced liver fibrosis. \emph{Hepatology} \textbf{74}, 2452--2466 (2021). \url{https://doi.org/10.1002/hep.31995}
\item
  Lefebvre, P. et al.~Interspecies NASH disease activity whole-genome profiling identifies a fibrogenic role of PPARα-regulated dermatopontin. \emph{JCI Insight} \textbf{2}, e92264 (2017). \url{https://doi.org/10.1172/jci.insight.92264}
\item
  Suppli, M. P. et al.~Hepatic transcriptome signatures in patients with varying degrees of nonalcoholic fatty liver disease compared with healthy normal-weight individuals. \emph{Am J Physiol Gastrointest Liver Physiol} \textbf{316}, G462--G472 (2019). \url{https://doi.org/10.1152/ajpgi.00358.2018}
\item
  Chen, T. et al.~Hepatocyte Smoothened activity controls susceptibility to insulin resistance and nonalcoholic fatty liver disease. \emph{Cell Mol Gastroenterol Hepatol} \textbf{15}, 949--970 (2023). \url{https://doi.org/10.1016/j.jcmgh.2022.12.008}
\item
  Fujiwara, N. et al.~Molecular signatures of long-term hepatocellular carcinoma risk in nonalcoholic fatty liver disease. \emph{Sci Transl Med} \textbf{14}, eabo4474 (2022). \url{https://doi.org/10.1126/scitranslmed.abo4474}
\item
  Fujiwara, N. et al.~AI-based phenotyping of hepatic fiber morphology to inform molecular alterations in metabolic dysfunction-associated steatotic liver disease. \emph{Hepatology} \textbf{83}, 575--590 (2026). \url{https://doi.org/10.1097/HEP.0000000000001360}
\item
  Zhao, Q. et al.~Global molecular landscape of early MASLD progression in human obesity. \emph{eLife} \textbf{14}, RP109534 (2026). \url{https://doi.org/10.7554/eLife.109534}
\item
  Powell, N. R. et al.~Clinically important alterations in pharmacogene expression in histologically severe nonalcoholic fatty liver disease. \emph{Nat Commun} \textbf{14}, 1474 (2023). \url{https://doi.org/10.1038/s41467-023-37209-1}
\item
  Verschuren, L. et al.~Development of a novel non-invasive biomarker panel for hepatic fibrosis in MASLD. \emph{Nat Commun} \textbf{15}, 4564 (2024). \url{https://doi.org/10.1038/s41467-024-48956-0}
\item
  Xanthakos, S. A. et al.~High prevalence of nonalcoholic fatty liver disease in adolescents undergoing bariatric surgery. \emph{Gastroenterology} \textbf{149}, 623--634.e8 (2015). \url{https://doi.org/10.1053/j.gastro.2015.05.039}
\item
  Yao, K. et al.~Transcriptomic profiling of a multiethnic pediatric NAFLD cohort reveals genes and pathways associated with disease. \emph{Hepatol Commun} \textbf{6}, 1598--1610 (2022). \url{https://doi.org/10.1002/hep4.1940}
\item
  Ahrens, M. et al.~DNA methylation analysis in nonalcoholic fatty liver disease suggests distinct disease-specific and remodeling signatures after bariatric surgery. \emph{Cell Metab} \textbf{18}, 296--302 (2013). \url{https://doi.org/10.1016/j.cmet.2013.07.004}
\item
  García-Nieto, E. et al.~Clinical and molecular characterization of steatotic liver disease in the setting of immune-mediated inflammatory diseases. \emph{JHEP Rep} \textbf{6}, 101167 (2024). \url{https://doi.org/10.1016/j.jhepr.2024.101167}
\item
  Yang, M. et al.~Integrative proteo-transcriptomic characterization of advanced fibrosis in metabolic dysfunction-associated steatotic liver disease. \emph{Cell Rep Med} \textbf{6}, 101935 (2025). \url{https://doi.org/10.1016/j.xcrm.2025.101935}
\item
  Steinhauser, M. L. et al.~The transcription factor ZNF469 regulates collagen production in liver fibrosis. \emph{JCI Insight} \textbf{10}, e182232 (2025). \url{https://doi.org/10.1172/jci.insight.182232}
\item
  Gerhard, G. S. et al.~Transcriptomic profiling of obesity-related nonalcoholic steatohepatitis (NASH). \emph{J Endocr Soc} \textbf{2}, 710--726 (2018). \url{https://doi.org/10.1210/js.2018-00122}
\item
  GTEx Consortium. The GTEx Consortium atlas of genetic regulatory effects across human tissues. \emph{Science} \textbf{369}, 1318--1330 (2020). \url{https://doi.org/10.1126/science.aaz1776}
\item
  Sugimoto, A. et al.~Hepatic stellate cells control liver zonation, size and functions via R-spondin 3. \emph{Nature} \textbf{640}, 752--761 (2025). \url{https://doi.org/10.1038/s41586-025-08677-w}
\item
  Alkhawaja, A. A. et al.~Liver single-nucleus multiome profiling reveals cell-type mechanisms for cardiometabolic traits. \emph{Am J Hum Genet} \textbf{113}, 133--148 (2026). \url{https://doi.org/10.1016/j.ajhg.2025.11.009}
\item
  Boel, F. et al.~Deep proteome profiling of metabolic dysfunction-associated steatotic liver disease. \emph{Commun Med} \textbf{5}, 56 (2025). \url{https://doi.org/10.1038/s43856-025-00780-3}
\item
  Tzouanas, C. N. et al.~Hepatic adaptation to chronic metabolic stress primes tumorigenesis. \emph{Cell} \textbf{189}, 435--460.e28 (2026). \url{https://doi.org/10.1016/j.cell.2025.11.031}
\item
  Verma, A. et al.~Diversity and scale: genetic architecture of 2068 traits in the VA Million Veteran Program. \emph{Science} \textbf{385}, eadj1182 (2024). \url{https://doi.org/10.1126/science.adj1182}
\item
  Kurki, M. I. et al.~FinnGen provides genetic insights from a well-phenotyped isolated population. \emph{Nature} \textbf{613}, 508--518 (2023). \url{https://doi.org/10.1038/s41586-022-05473-8}
\item
  Karczewski, K. J. et al.~Pan-UK Biobank genome-wide association analyses enhance discovery and resolution of ancestry-enriched effects. \emph{Nat Genet} \textbf{57}, 2408--2417 (2025). \url{https://doi.org/10.1038/s41588-025-02335-7}
\item
  Ghouse, J. et al.~Integrative common and rare variant analyses provide insights into the genetic architecture of liver cirrhosis. \emph{Nat Genet} \textbf{56}, 827--837 (2024). \url{https://doi.org/10.1038/s41588-024-01720-y}
\item
  Chen, Y. et al.~Genomic atlas of the plasma metabolome prioritizes metabolites implicated in human diseases. \emph{Nat Genet} \textbf{55}, 44--53 (2023). \url{https://doi.org/10.1038/s41588-022-01270-1}
\item
  Ottensmann, L. et al.~Genome-wide association analysis of plasma lipidome identifies 495 genetic associations. \emph{Nat Commun} \textbf{14}, 6934 (2023). \url{https://doi.org/10.1038/s41467-023-42532-8}
\item
  Zhao, J. H. et al.~Genetics of circulating inflammatory proteins identifies drivers of immune-mediated disease risk and therapeutic targets. \emph{Nat Immunol} \textbf{24}, 1540--1551 (2023). \url{https://doi.org/10.1038/s41590-023-01588-w}
\item
  Du, K. et al.~Targeting senescent hepatocytes for treatment of metabolic dysfunction-associated steatotic liver disease and multi-organ dysfunction. \emph{Nat Commun} \textbf{16}, 3038 (2025). \url{https://doi.org/10.1038/s41467-025-57616-w}
\item
  Qiu, X. et al.~THEMIS attenuates MASH by suppressing disease-associated hepatocyte induction and hepatocyte senescence in mice. \emph{J Clin Invest} \textbf{136}, e199303 (2026). \url{https://doi.org/10.1172/JCI199303}
\item
  Fisher, C. D. et al.~Hepatic cytochrome P450 enzyme alterations in humans with progressive stages of nonalcoholic fatty liver disease. \emph{Drug Metab Dispos} \textbf{37}, 2087--2094 (2009). \url{https://doi.org/10.1124/dmd.109.027466}
\item
  Murphy, S. K. et al.~Relationship between methylome and transcriptome in patients with nonalcoholic fatty liver disease. \emph{Gastroenterology} \textbf{145}, 1076--1087 (2013). \url{https://doi.org/10.1053/j.gastro.2013.07.047}
\item
  Zubiete-Franco, I. et al.~Methionine and S-adenosylmethionine levels are critical regulators of PP2A activity modulating lipophagy during steatosis. \emph{J Hepatol} \textbf{64}, 409--418 (2016). \url{https://doi.org/10.1016/j.jhep.2015.08.037}
\item
  Fernández-Tussy, P. et al.~miR-873-5p targets mitochondrial GNMT-Complex II interface contributing to non-alcoholic fatty liver disease. \emph{Mol Metab} \textbf{29}, 40--54 (2019). \url{https://doi.org/10.1016/j.molmet.2019.08.008}
\item
  Männistö, V. et al.~Total liver phosphatidylcholine content associates with non-alcoholic steatohepatitis and glycine N-methyltransferase expression. \emph{Liver Int} \textbf{39}, 1895--1905 (2019). \url{https://doi.org/10.1111/liv.14174}
\item
  Gaidatzis, D., Burger, L., Florescu, M. \& Stadler, M. B. Analysis of intronic and exonic reads in RNA-seq data characterizes transcriptional and post-transcriptional regulation. \emph{Nat Biotechnol} \textbf{33}, 722--729 (2015). \url{https://doi.org/10.1038/nbt.3269}
\item
  1000 Genomes Project Consortium. A global reference for human genetic variation. \emph{Nature} \textbf{526}, 68--74 (2015). \url{https://doi.org/10.1038/nature15393}
\item
  Frankish, A. et al.~GENCODE: reference annotation for the human and mouse genomes in 2023. \emph{Nucleic Acids Res} \textbf{51}, D942--D949 (2023). \url{https://doi.org/10.1093/nar/gkac1071}
\item
  Chen, S. et al.~A genomic mutational constraint map using variation in 76,156 human genomes. \emph{Nature} \textbf{625}, 92--100 (2024). \url{https://doi.org/10.1038/s41586-023-06045-0}
\item
  Tsherniak, A. et al.~Defining a cancer dependency map. \emph{Cell} \textbf{170}, 564--576.e16 (2017). \url{https://doi.org/10.1016/j.cell.2017.06.010}
\item
  Uhlén, M. et al.~Tissue-based map of the human proteome. \emph{Science} \textbf{347}, 1260419 (2015). \url{https://doi.org/10.1126/science.1260419}
\item
  Hart, T. et al.~Evaluation and design of genome-wide CRISPR/SpCas9 knockout screens. \emph{G3} \textbf{7}, 2719--2727 (2017). \url{https://doi.org/10.1534/g3.117.041277}
\item
  Garcia-Alonso, L., Holland, C. H., Ibrahim, M. M., Turei, D. \& Saez-Rodriguez, J. Benchmark and integration of resources for the estimation of human transcription factor activities. \emph{Genome Res} \textbf{29}, 1363--1375 (2019). \url{https://doi.org/10.1101/gr.240663.118}
\item
  Liberzon, A. et al.~The Molecular Signatures Database hallmark gene set collection. \emph{Cell Syst} \textbf{1}, 417--425 (2015). \url{https://doi.org/10.1016/j.cels.2015.12.004}
\item
  Dobin, A. et al.~STAR: ultrafast universal RNA-seq aligner. \emph{Bioinformatics} \textbf{29}, 15--21 (2013). \url{https://doi.org/10.1093/bioinformatics/bts635}
\item
  Liao, Y., Smyth, G. K. \& Shi, W. featureCounts: an efficient general purpose program for assigning sequence reads to genomic features. \emph{Bioinformatics} \textbf{30}, 923--930 (2014). \url{https://doi.org/10.1093/bioinformatics/btt656}
\item
  Robinson, M. D., McCarthy, D. J. \& Smyth, G. K. edgeR: a Bioconductor package for differential expression analysis of digital gene expression data. \emph{Bioinformatics} \textbf{26}, 139--140 (2010). \url{https://doi.org/10.1093/bioinformatics/btp616}
\item
  Ritchie, M. E. et al.~limma powers differential expression analyses for RNA-sequencing and microarray studies. \emph{Nucleic Acids Res} \textbf{43}, e47 (2015). \url{https://doi.org/10.1093/nar/gkv007}
\item
  Law, C. W., Chen, Y., Shi, W. \& Smyth, G. K. voom: precision weights unlock linear model analysis tools for RNA-seq read counts. \emph{Genome Biol} \textbf{15}, R29 (2014). \url{https://doi.org/10.1186/gb-2014-15-2-r29}
\item
  Liu, R. et al.~Why weight? Modelling sample and observational level variability improves power in RNA-seq analyses. \emph{Nucleic Acids Res} \textbf{43}, e97 (2015). \url{https://doi.org/10.1093/nar/gkv412}
\item
  Stephens, M. False discovery rates: a new deal. \emph{Biostatistics} \textbf{18}, 275--294 (2017). \url{https://doi.org/10.1093/biostatistics/kxw041}
\item
  Wang, G., Sarkar, A., Carbonetto, P. \& Stephens, M. A simple new approach to variable selection in regression, with application to genetic fine mapping. \emph{J R Stat Soc Series B Stat Methodol} \textbf{82}, 1273--1300 (2020). \url{https://doi.org/10.1111/rssb.12388}
\item
  Giambartolomei, C. et al.~Bayesian test for colocalisation between pairs of genetic association studies using summary statistics. \emph{PLoS Genet} \textbf{10}, e1004383 (2014). \url{https://doi.org/10.1371/journal.pgen.1004383}
\item
  Wallace, C. A more accurate method for colocalisation analysis allowing for multiple causal variants. \emph{PLoS Genet} \textbf{17}, e1009440 (2021). \url{https://doi.org/10.1371/journal.pgen.1009440}
\item
  Weissbrod, O. et al.~Functionally informed fine-mapping and polygenic localization of complex trait heritability. \emph{Nat Genet} \textbf{52}, 1355--1363 (2020). \url{https://doi.org/10.1038/s41588-020-00735-5}
\item
  Dai, Q. et al.~OTTERS: a powerful TWAS framework leveraging summary-level reference data. \emph{Nat Commun} \textbf{14}, 1271 (2023). \url{https://doi.org/10.1038/s41467-023-36862-w}
\item
  Yuan, K. et al.~Fine-mapping across diverse ancestries drives the discovery of putative causal variants underlying human complex traits and diseases. \emph{Nat Genet} \textbf{56}, 1841--1850 (2024). \url{https://doi.org/10.1038/s41588-024-01870-z}
\item
  Gao, B. \& Zhou, X. MESuSiE enables scalable and powerful multi-ancestry fine-mapping of causal variants in genome-wide association studies. \emph{Nat Genet} \textbf{56}, 170--179 (2024). \url{https://doi.org/10.1038/s41588-023-01604-7}
\item
  McLaren, W. et al.~The Ensembl Variant Effect Predictor. \emph{Genome Biol} \textbf{17}, 122 (2016). \url{https://doi.org/10.1186/s13059-016-0974-4}
\item
  Hinrichs, A. S. et al.~The UCSC Genome Browser Database: update 2006. \emph{Nucleic Acids Res} \textbf{34}, D590--D598 (2006). \url{https://doi.org/10.1093/nar/gkj144}
\item
  Lawrence, M., Gentleman, R. \& Carey, V. rtracklayer: an R package for interfacing with genome browsers. \emph{Bioinformatics} \textbf{25}, 1841--1842 (2009). \url{https://doi.org/10.1093/bioinformatics/btp328}
\item
  Oh, S. et al.~HGNChelper: identification and correction of invalid gene symbols for human and mouse. \emph{F1000Res} \textbf{9}, 1493 (2020). \url{https://doi.org/10.12688/f1000research.28033.1}
\item
  DeTomaso, D. \& Yosef, N. Hotspot identifies informative gene modules across modalities of single-cell genomics. \emph{Cell Syst} \textbf{12}, 446--456.e9 (2021). \url{https://doi.org/10.1016/j.cels.2021.04.005}
\item
  Lopez, R., Regier, J., Cole, M. B., Jordan, M. I. \& Yosef, N. Deep generative modeling for single-cell transcriptomics. \emph{Nat Methods} \textbf{15}, 1053--1058 (2018). \url{https://doi.org/10.1038/s41592-018-0229-2}
\item
  Yang, S. et al.~Decontamination of ambient RNA in single-cell RNA-seq with DecontX. \emph{Genome Biol} \textbf{21}, 57 (2020). \url{https://doi.org/10.1186/s13059-020-1950-6}
\item
  Wolock, S. L., Lopez, R. \& Klein, A. M. Scrublet: computational identification of cell doublets in single-cell transcriptomic data. \emph{Cell Syst} \textbf{8}, 281--291.e9 (2019). \url{https://doi.org/10.1016/j.cels.2018.11.005}
\item
  Dimitrov, D. et al.~Comparison of methods and resources for cell-cell communication inference from single-cell RNA-Seq data. \emph{Nat Commun} \textbf{13}, 3224 (2022). \url{https://doi.org/10.1038/s41467-022-30755-0}
\item
  Badia-i-Mompel, P. et al.~decoupleR: ensemble of computational methods to infer biological activities from omics data. \emph{Bioinform Adv} \textbf{2}, vbac016 (2022). \url{https://doi.org/10.1093/bioadv/vbac016}
\item
  Chu, T., Wang, Z., Pe'er, D. \& Danko, C. G. Cell type and gene expression deconvolution with BayesPrism enables Bayesian integrative analysis across bulk and single-cell RNA sequencing in oncology. \emph{Nat Cancer} \textbf{3}, 505--517 (2022). \url{https://doi.org/10.1038/s43018-022-00356-3}
\item
  Hu, M. \& Chikina, M. InstaPrism: an R package for fast implementation of BayesPrism. \emph{Bioinformatics} \textbf{40}, btae440 (2024). \url{https://doi.org/10.1093/bioinformatics/btae440}
\item
  Kleshchevnikov, V. et al.~Cell2location maps fine-grained cell types in spatial transcriptomics. \emph{Nat Biotechnol} \textbf{40}, 661--671 (2022). \url{https://doi.org/10.1038/s41587-021-01139-4}
\item
  Traag, V. A., Waltman, L. \& van Eck, N. J. From Louvain to Leiden: guaranteeing well-connected communities. \emph{Sci Rep} \textbf{9}, 5233 (2019). \url{https://doi.org/10.1038/s41598-019-41695-z}
\item
  McInnes, L., Healy, J., Saul, N. \& Großberger, L. UMAP: uniform manifold approximation and projection. \emph{J Open Source Softw} \textbf{3}, 861 (2018). \url{https://doi.org/10.21105/joss.00861}
\item
  Hu, M. et al.~ScaleSC: a superfast and scalable single-cell RNA-seq analysis tool. \emph{Bioinform Adv} \textbf{5}, vbaf167 (2024). \url{https://doi.org/10.1093/bioadv/vbaf167}
\item
  Korsunsky, I. et al.~Fast, sensitive and accurate integration of single-cell data with Harmony. \emph{Nat Methods} \textbf{16}, 1289--1296 (2019). \url{https://doi.org/10.1038/s41592-019-0619-0}
\item
  Domínguez Conde, C. et al.~Cross-tissue immune cell analysis reveals tissue-specific features in humans. \emph{Science} \textbf{376}, eabl5197 (2022). \url{https://doi.org/10.1126/science.abl5197}
\item
  Chang, C. C. et al.~Second-generation PLINK: rising to the challenge of larger and richer datasets. \emph{GigaScience} \textbf{4}, s13742-015-0047-8 (2015). \url{https://doi.org/10.1186/s13742-015-0047-8}
\item
  Zhang, K. et al.~A fast, scalable and versatile tool for analysis of single-cell omics data. \emph{Nat Methods} \textbf{21}, 217--227 (2024). \url{https://doi.org/10.1038/s41592-023-02139-9}
\item
  Zhang, Y. et al.~Model-based analysis of ChIP-Seq (MACS). \emph{Genome Biol} \textbf{9}, R137 (2008). \url{https://doi.org/10.1186/gb-2008-9-9-r137}
\item
  Grotzinger, A. D. et al.~Genomic structural equation modelling provides insights into the multivariate genetic architecture of complex traits. \emph{Nat Hum Behav} \textbf{3}, 513--525 (2019). \url{https://doi.org/10.1038/s41562-019-0566-x}
\item
  Pearce, J. D. et al.~TranscriptFormer: a generative cell atlas across 1.5 billion years of evolution. \emph{Science} \textbf{393}, aec8514 (2026). \url{https://doi.org/10.1126/science.aec8514}
\item
  Chen, H. \& Theodoris, C. V. Quantized multi-task learning for context-specific representations of gene network dynamics. Preprint at bioRxiv \url{https://doi.org/10.1101/2024.08.16.608180} (2024).
\item
  Kang, B. et al.~A large-scale foundation model for bulk transcriptomes. Preprint at bioRxiv \url{https://doi.org/10.1101/2025.06.11.659222} (2025).
\item
  Nguyen, E. et al.~HyenaDNA: long-range genomic sequence modeling at single nucleotide resolution. \emph{Adv Neural Inf Process Syst} \textbf{36}, 43177--43201 (2023). \url{https://arxiv.org/abs/2306.15794}
\item
  Dalla-Torre, H. et al.~Nucleotide Transformer: building and evaluating robust foundation models for human genomics. \emph{Nat Methods} \textbf{22}, 287--297 (2025). \url{https://doi.org/10.1038/s41592-024-02523-z}
\item
  Rosen, Y. et al.~Universal Cell Embeddings: a foundation model for cell biology. Preprint at bioRxiv \url{https://doi.org/10.1101/2023.11.28.568918} (2023).
\item
  Heimberg, G. et al.~A cell atlas foundation model for scalable search of similar human cells. \emph{Nature} \textbf{638}, 1085--1094 (2025). \url{https://doi.org/10.1038/s41586-024-08411-y}
\item
  Freedman, D. \& Lane, D. A nonstochastic interpretation of reported significance levels. \emph{J Bus Econ Stat} \textbf{1}, 292--298 (1983). \url{https://doi.org/10.1080/07350015.1983.10509354}
\item
  Long, J. S. \& Ervin, L. H. Using heteroscedasticity consistent standard errors in the linear regression model. \emph{Am Stat} \textbf{54}, 217--224 (2000). \url{https://doi.org/10.1080/00031305.2000.10474549}
\item
  Knapp, G. \& Hartung, J. Improved tests for a random effects meta-regression with a single covariate. \emph{Stat Med} \textbf{22}, 2693--2710 (2003). \url{https://doi.org/10.1002/sim.1482}
\item
  Mantel, N. \& Haenszel, W. Statistical aspects of the analysis of data from retrospective studies of disease. \emph{J Natl Cancer Inst} \textbf{22}, 719--748 (1959). \url{https://doi.org/10.1093/jnci/22.4.719}
\item
  Robins, J., Breslow, N. \& Greenland, S. Estimators of the Mantel-Haenszel variance consistent in both sparse data and large-strata limiting models. \emph{Biometrics} \textbf{42}, 311--323 (1986). \url{https://doi.org/10.2307/2531052}
\item
  Moran, P. A. P. Notes on continuous stochastic phenomena. \emph{Biometrika} \textbf{37}, 17--23 (1950). \url{https://doi.org/10.1093/biomet/37.1-2.17}
\end{enumerate}
